%Daten zur Institution

%\input{Dozentdaten}

%\renewcommand{\fachbereich}{Fachbereich}

%\renewcommand{\dozent}{Prof. Dr. . }

%Klausurdaten

\renewcommand{\klausurgebiet}{ }

\renewcommand{\klausurtyp}{ }

\renewcommand{\klausurdatum}{. 20}

\klausurvorspann {\fachbereich} {\klausurdatum} {\dozent} {\klausurgebiet} {\klausurtyp}

%Daten für folgende Punktetabelle

\renewcommand{\aeins}{  3  }

\renewcommand{\azwei}{  3   }

\renewcommand{\adrei}{  6  }

\renewcommand{\avier}{  5  }

\renewcommand{\afuenf}{  5  }

\renewcommand{\asechs}{  2  }

\renewcommand{\asieben}{  5  }

\renewcommand{\aacht}{  5  }

\renewcommand{\aneun}{  7  }

\renewcommand{\azehn}{  3  }

\renewcommand{\aelf}{  2  }

\renewcommand{\azwoelf}{  5  }

\renewcommand{\adreizehn}{  9   }

\renewcommand{\avierzehn}{  5  }

\renewcommand{\afuenfzehn}{  65  }

\renewcommand{\asechzehn}{    }

\renewcommand{\asiebzehn}{    }

\renewcommand{\aachtzehn}{    }

\renewcommand{\aneunzehn}{    }

\renewcommand{\azwanzig}{    }

\renewcommand{\aeinundzwanzig}{    }

\renewcommand{\azweiundzwanzig}{    }

\renewcommand{\adreiundzwanzig}{    }

\renewcommand{\avierundzwanzig}{    }

\renewcommand{\afuenfundzwanzig}{    }

\renewcommand{\asechsundzwanzig}{    }

\punktetabellevierzehn

\klausurnote

\newpage

\setcounter{section}{0}

__NOEDITSECTION__

\inputaufgabeklausurloesung
{3}

{

Definiere die folgenden
\zusatzklammer {kursiv gedruckten} {} {} Begriffe.
\aufzaehlungsechs{Der
\stichwort {Krümmungskreis} {}
zu einer zweifach
\definitionsverweis {stetig differenzierbare}{}{}
\definitionsverweis {bogenparametrisierte}{}{}
\definitionsverweis {Kurve}{}{}
Es sei
\maabbdisp {\gamma} {I} {\R^2
} {}
in einem Punkt

\mavergleichskette
{\vergleichskette
{ t_0
}
{ \in }{ I
}
{  }{
}
{    }{
}
{   }{
}
}
{}{}{.}

}{Eine \stichwort {abgeschlossene Untermannigfaltigkeit} {}
\mathl{M \subseteq N}{} einer differenzierbaren Mannigfaltigkeit $N$.

}{Der \stichwort {Tangentialraum} {} in einem Punkt
\mathl{P\in M}{} einer differenzierbaren Mannigfaltigkeit $M$.

}{Eine
\stichwort {orientierte} {}
\definitionsverweis {differenzierbare Mannigfaltigkeit}{}{}
$M$.

}{Der
\stichwort {Rand} {}
einer
\definitionsverweis {Mannigfaltigkeit mit Rand}{}{.}

}{Die
\stichwort {Christoffelsymbole} {}

\mathl{\Gamma_{ij}^k}{} für den Levi-Civita-Zusammenhang zu einer riemannschen Struktur  $\left(  g_{ij}  \right)$ auf einer offenen Menge

\mavergleichskette
{\vergleichskette
{ U
}
{ \subseteq }{ \R^n
}
{  }{
}
{    }{
}
{   }{
}
}
{}{}{.}
}

}
{

\aufzaehlungsechs{Man nennt den Kreis mit dem Radius

\mavergleichskettedisp
{\vergleichskette
{ r
}
{ =} {  { \frac{  1 }{ \betrag {  \gamma_1'(t_0) \gamma_2^{\prime \prime} (t_0) - \gamma_2'(t_0) \gamma_1^{\prime \prime} (t_0)  }  } }
}
{ } {
}
{ } {
}
{ } {
}
}
{}{}{}
und dem Mittelpunkt

\mavergleichskettedisp
{\vergleichskette
{ M
}
{ =} { \gamma(t_0) +  { \frac{  1 }{  \gamma_1'(t_0) \gamma_2^{\prime \prime} (t_0) - \gamma_2'(t_0) \gamma_1^{\prime \prime} (t_0)  } }  \begin{pmatrix} - \gamma_2'(t_0) \\ \gamma_1'(t_0)   \end{pmatrix}
}
{ } {
}
{ } {
}
{ } {
}
}
{}{}{}
den
Krümmungskreis
zu $\gamma$ in $t_0$.
}{Eine abgeschlossene Teilmenge
\mathl{M \subseteq N}{} heißt abgeschlossene Untermannigfaltigkeit, wenn es zu jedem Punkt
\mathl{P \in M}{} eine
\definitionsverweis {Karte}{}{}
gibt mit
\mathl{P \in W \subseteq N}{}
\definitionsverweis {offen}{}{,}
\maabb {\theta} { W } { W'
} {,}

\mathl{W' \subseteq \R^n}{} offen und mit

\mavergleichskettedisp
{\vergleichskette
{  M \cap W
}
{ =} {\theta^{-1} (( \R^m  \times\{0\})  \cap W' )
}
{ } {
}
{ } {
}
{ } {
}
}
{}{}{.}
}{Der Tangentialraum
\mathl{T_PM}{} besteht aus allen Äquivalenzklassen von tangential äquivalenten differenzierbaren Wegen durch diesen Punkt.
}{Eine
\definitionsverweis {differenzierbare Mannigfaltigkeit}{}{}
$M$ mit einem
\definitionsverweis {Atlas}{}{}

\mathl{(U_i, V_i, \alpha_i)}{} heißt orientiert, wenn jede Karte
\definitionsverweis {orientiert}{}{}
ist und wenn sämtliche Kartenwechsel
\definitionsverweis {orientierungstreu}{}{}
sind.
}{Der Rand von $M$ ist durch

\mavergleichskettedisp
{\vergleichskette
{ \partial M
}
{ =} {  \left\{ x \in M  \mid   \alpha_i(x) \in  \partial H \cap V_i  \text{ für } \text{ein } i \in I         \right\}
}
{ } {
}
{ } {
}
{ } {
}
}
{}{}{}
definiert, wobei  $\alpha_i$ Karten sind.
}{Es sei
\mathl{\left(  g^{k l}  \right)}{} die inverse Matrix zu
\mathl{\left(  g_{ij}  \right)}{.} Man nennt
\zusatzklammer {die auf $U$ definierten reellwertigen Funktionen} {} {}

\mavergleichskettedisp
{\vergleichskette
{\Gamma_{ij}^k
}
{ =} {  { \frac{  1 }{ 2 } } \sum_{l = 1}^n g^{kl}(\partial_ig_{jl}+\partial_jg_{il}-\partial_lg_{ij})
}
{ } {
}
{ } {
}
{ } {
}
}
{}{}{}
die
Christoffelsymbole
für den Levi-Civita-Zusammenhang.
}

 }

__NOEDITSECTION__

\inputaufgabeklausurloesung
{3}

{

Formuliere die folgenden Sätze.
\aufzaehlungdrei{Der Satz über das Bild der
\definitionsverweis {Gauß-Abbildung}{}{}
eine Hyperfläche.}{Die Formel für den Flächeninhalt einer Rotationsfläche zu einer
\definitionsverweis {differenzierbaren Kurve}{}{}
\maabbeledisp {\gamma} {]a,b[} {\R^2
} { t } {(x(t),y(t))
} {,}
mit

\mavergleichskette
{\vergleichskette
{  y(t)
}
{ > }{  0
}
{  }{
}
{    }{
}
{   }{
}
}
{}{}{.}}{Die
\stichwort {Produktregel} {}
für die äußere Ableitung.}

}
{

\aufzaehlungdrei{\faktsituation {Es sei

\mavergleichskette
{\vergleichskette
{ W
}
{ \subseteq }{\R^n
}
{  }{
}
{    }{
}
{   }{
}
}
{}{}{}
eine
\definitionsverweis {offene Teilmenge}{}{}
und sei
\maabbdisp {h} { W } {\R
} {}
eine
\definitionsverweis {stetig differenzierbare Funktion}{}{.}}
\faktvoraussetzung {Die
\definitionsverweis {Faser}{}{}

\mavergleichskette
{\vergleichskette
{ Y
}
{ = }{ h^{-1}(c)
}
{  }{
}
{    }{
}
{   }{
}
}
{}{}{}
von $h$ zu

\mavergleichskette
{\vergleichskette
{ c
}
{ \in }{\R
}
{  }{
}
{    }{
}
{   }{
}
}
{}{}{}
sei
\definitionsverweis {kompakt}{}{}
und in jedem Punkt
\definitionsverweis {regulär}{}{.}}
\faktfolgerung {Dann ist die
\definitionsverweis {Gauß-Abbildung}{}{}
\maabbdisp {} { Y } { S^{n-1}
} {}
surjektiv.}
\faktzusatz {}
\faktzusatz {}}{\faktsituation {Es sei
\maabbeledisp {\gamma} {]a,b[} {\R^2
} { t } {(x(t),y(t))
} {,}
eine
\definitionsverweis {differenzierbare Kurve}{}{}
mit

\mavergleichskette
{\vergleichskette
{ y(t)
}
{ > }{ 0
}
{  }{
}
{    }{
}
{   }{
}
}
{}{}{,}}
\faktvoraussetzung {die einen
\definitionsverweis {Diffeomorphismus}{}{}
zu

\mavergleichskette
{\vergleichskette
{M
}
{ = }{ \gamma(I)
}
{  }{
}
{    }{
}
{   }{
}
}
{}{}{}
induziere, wobei

\mavergleichskette
{\vergleichskette
{M
}
{ \subseteq }{ G
}
{  }{
}
{    }{
}
{   }{
}
}
{}{}{}
eine eindimensionale
\definitionsverweis {abgeschlossene Untermannigfaltigkeit}{}{}
in einer offenen Menge

\mavergleichskette
{\vergleichskette
{G
}
{ \subseteq }{\R \times \R_+
}
{  }{
}
{    }{
}
{   }{
}
}
{}{}{}
sei.}
\faktfolgerung {Dann ist die zugehörige
\definitionsverweis {Rotationsfläche}{}{}
eine abgeschlossene Untermannigfaltigkeit von
\mathl{\R^3}{} ohne die $x$-Achse, und ihr Flächeninhalt ist gleich

\mathdisp {2 \pi \int_{  a  }^{  b  } \sqrt{ (x'(t))^2+ (y'(t))^2 } \cdot y(t) \, d t} { . }
}
\faktzusatz {}
\faktzusatz {}}{\faktsituation {Es sei

\mavergleichskette
{\vergleichskette
{ U
}
{ \subseteq }{ \R^n
}
{  }{
}
{    }{
}
{   }{
}
}
{}{}{}
\definitionsverweis {offen}{}{}
und es seien

\mavergleichskette
{\vergleichskette
{ \omega
}
{ \in }{
{ \mathcal E }^{  k  }_{ 1 } (  U   )
}
{  }{
}
{    }{
}
{   }{
}
}
{}{}{}
und

\mavergleichskette
{\vergleichskette
{ \tau
}
{ \in }{
{ \mathcal E }^{  \ell }_{ 1 } (  U   )
}
{  }{
}
{    }{
}
{   }{
}
}
{}{}{}
differenzierbare Differentialformen auf $U$.}
\faktfolgerung {Dann gilt die Produktregel

\mavergleichskettedisp
{\vergleichskette
{ d( \omega \wedge \tau)
}
{ =} { ( d \omega) \wedge \tau  + (-1)^k \omega  \wedge  (d \tau)
}
{ } {
}
{ } {
}
{ } {
}
}
{}{}{.}}
\faktzusatz {}
\faktzusatz {}}

 }

__NOEDITSECTION__

\inputaufgabeklausurloesung
{6 (1+2+3)}

{

Es sei $Y$ die Hyperfläche, die durch die Bedingung

\mavergleichskettedisp
{\vergleichskette
{ 2x^2 +3y^2+5z^2
}
{ =} { 10
}
{ } {
}
{ } {
}
{ } {
}
}
{}{}{}
im $\R^3$ gegeben ist. Es sei

\mavergleichskette
{\vergleichskette
{ P
}
{ = }{ \begin{pmatrix}  -1 \\ -1\\  1   \end{pmatrix}
}
{ \in }{ Y
}
{    }{
}
{   }{
}
}
{}{}{}
und

\mavergleichskette
{\vergleichskette
{ w
}
{ = }{ \begin{pmatrix}  -3 \\ 2 \\  4   \end{pmatrix}
}
{ \in }{ \R^3
}
{    }{
}
{   }{
}
}
{}{}{.}
\aufzaehlungdrei{Bestimme den
\definitionsverweis {Tangentialraum}{}{}

\mathl{T_PY}{.}
}{Bestimme die
\definitionsverweis {orthogonale Zerlegung}{}{}
$w$ in die tangentiale und in die orthogonale Komponente zum Punkt $P$.
}{Realisiere die tangentiale Komponente von $w$ in $T_PY$ durch eine differenzierbare Kurve, die ganz auf $Y$  verläuft.
}

}
{

\aufzaehlungdrei{Das totale Differential ist

\mathdisp {\left(  4x  , \,   6y  , \,   10z                 \right)} { , }

in $P$ führt das auf

\mavergleichskettedisp
{\vergleichskette
{ T_PY
}
{ =} { \left\{  \begin{pmatrix}  a  \\ b \\ c   \end{pmatrix}   \mid   -4a+6b+10c = 0         \right\}
}
{ =} { \left\{  \begin{pmatrix}  a  \\ b \\ c   \end{pmatrix}   \mid      \left\langle  \begin{pmatrix}  a  \\ b \\ c   \end{pmatrix}   , \begin{pmatrix}  -4 \\ -6\\  10   \end{pmatrix}     \right\rangle = 0         \right\}
}
{ } {
}
{ } {
}
}
{}{}{.}
}{Wir machen den Ansatz

\mavergleichskettedisp
{\vergleichskette
{  \begin{pmatrix}  -3 \\ 2\\  4   \end{pmatrix}
}
{ =} { s  \begin{pmatrix}  -4 \\ -6\\  10   \end{pmatrix} + v
}
{ } {
}
{ } {
}
{ } {
}
}
{}{}{}
mit einem gesuchten Skalar

\mavergleichskette
{\vergleichskette
{ s
}
{ \in }{ \R
}
{  }{
}
{    }{
}
{   }{
}
}
{}{}{}
und dem gesuchten tangentialen Vektor $v$. Es ist

\mavergleichskettedisp
{\vergleichskette
{ \left\langle  \begin{pmatrix}  -3 \\ 2\\  4   \end{pmatrix}  ,  \begin{pmatrix}  -4 \\ -6\\  10   \end{pmatrix}   \right\rangle
}
{ =} { 12-12+40
}
{ =} { 40
}
{ } {
}
{ } {
}
}
{}{}{.}
Wegen

\mavergleichskettedisp
{\vergleichskette
{ \left\langle  \begin{pmatrix}  -4 \\ -6\\  10   \end{pmatrix}  ,  \begin{pmatrix}  -4 \\ -6\\  10   \end{pmatrix}   \right\rangle
}
{ =} { 16 +36 + 100
}
{ =} { 152
}
{ } {
}
{ } {
}
}
{}{}{}
ist

\mavergleichskettedisp
{\vergleichskette
{ s
}
{ =} { { \frac{   40 }{  152 } }
}
{ =} { { \frac{   5 }{  19 } }
}
{ } {
}
{ } {
}
}
{}{}{}
und damit

\mavergleichskettedisp
{\vergleichskette
{ v
}
{ =} { \begin{pmatrix}  -3 \\ 2\\  4   \end{pmatrix}   - { \frac{   5 }{  19 } }  \begin{pmatrix}  -4 \\ -6\\  10   \end{pmatrix}
}
{ =} { { \frac{   1  }{  19 } }  \begin{pmatrix}  -57+20 \\ 38 +30\\  76 -50   \end{pmatrix}
}
{ =} { { \frac{   1  }{  19 } }  \begin{pmatrix}  -37 \\ 68\\  26   \end{pmatrix}
}
{ } {
}
}
{}{}{.}
Die orthogonale Zerlegung ist also

\mavergleichskettedisp
{\vergleichskette
{  w
}
{ =} { \begin{pmatrix}  -3 \\ 2\\  4   \end{pmatrix}
}
{ =} { { \frac{   5 }{  19 } }  \begin{pmatrix}  -4 \\ -6\\  10   \end{pmatrix} + { \frac{   1  }{  19 } }  \begin{pmatrix}  -37 \\ 68\\  26   \end{pmatrix}
}
{ } {
}
{ } {
}
}
{}{}{.}
}{Wir schreiben die Flächengleichung

\mavergleichskette
{\vergleichskette
{ 2x^2 +3y^2+5z^2
}
{ = }{ 10
}
{  }{
}
{    }{
}
{   }{
}
}
{}{}{}
als

\mavergleichskettedisp
{\vergleichskette
{ z
}
{ =} { \pm \sqrt{ 10- 2x^2 - 3y^2 }
}
{ } {
}
{ } {
}
{ } {
}
}
{}{}{,}
wobei im Punkt $P$ die positive Wurzel zu nehmen ist. Dies ist eine Parametrisierung der Fläche in einer offenen Umgebung von $P$, wobei der Tangentialraum in
\mathl{\begin{pmatrix}  -1 \\ -1   \end{pmatrix}}{} in den Tangentialraum
\mathl{T_PY}{} abgebildet wird.
}

 }

__NOEDITSECTION__

\inputaufgabeklausurloesung
{5}

{

Es sei
\maabbeledisp {\gamma} {\R} { \R^2
} { t } { \begin{pmatrix}  \cos \left(  \omega t \right)  \\ \sin  \left(  \omega t \right)    \end{pmatrix}
} {,}
mit einem

\mavergleichskette
{\vergleichskette
{ \omega
}
{ > }{ 0
}
{  }{
}
{    }{
}
{   }{
}
}
{}{}{.}
Man zeige, dass es keine andere Kreisbewegung mit konstanter Geschwindigkeitsnorm gibt, die mit $\gamma$ für

\mavergleichskette
{\vergleichskette
{ t
}
{ = }{ 0
}
{  }{
}
{    }{
}
{   }{
}
}
{}{}{}
bis zur zweiten Ableitung übereinstimmt.

}
{

Es ist

\mavergleichskettedisp
{\vergleichskette
{ \gamma(0)
}
{ =} { \begin{pmatrix}  1  \\ 0   \end{pmatrix}
}
{ } {
}
{ } {
}
{ } {
}
}
{}{}{,}

\mavergleichskettedisp
{\vergleichskette
{ \gamma'(0)
}
{ =} { \begin{pmatrix}  0  \\ \omega   \end{pmatrix}
}
{ } {
}
{ } {
}
{ } {
}
}
{}{}{}
und

\mavergleichskettedisp
{\vergleichskette
{ \gamma^{\prime \prime} (0)
}
{ =} { \begin{pmatrix}  - \omega^2  \\ 0   \end{pmatrix}
}
{ } {
}
{ } {
}
{ } {
}
}
{}{}{.}
Eine Kreisbewegung, die mit der vorgegebenen Bewegung für

\mavergleichskette
{\vergleichskette
{ t
}
{ = }{ 0
}
{  }{
}
{    }{
}
{   }{
}
}
{}{}{}
bis zur zweiten Ableitung übereinstimmt, muss insbesondere die gleiche Tangente in diesem Punkt haben. Daher muss der Kreismittelpunkt auf der $x$-Achse liegen. Ferner muss er links von
\mathl{\begin{pmatrix}  1  \\ 0   \end{pmatrix}}{} liegen, da die Beschleunigung nach links wirkt. Wir können daher eine gesuchte uniforme Kreisbewegung mit dem Mittelpunkt
\mathl{\begin{pmatrix}  x  \\ 0   \end{pmatrix}}{}
\zusatzklammer {
\mavergleichskettek
{\vergleichskettek
{ x
}
{ < }{  1
}
{  }{
}
{    }{
}
{   }{
}
}
{}{}{}} {} {}
und dem Radius
\mathl{1-x}{}  als

\mavergleichskettedisp
{\vergleichskette
{  \delta (t)
}
{ =} {  \begin{pmatrix}  x  \\ 0   \end{pmatrix} + (1-x) \begin{pmatrix}  \cos \left(  u t \right)  \\ \sin  \left(  u t \right)     \end{pmatrix}
}
{ } {
}
{ } {
}
{ } {
}
}
{}{}{}
ansetzen. Es ist

\mavergleichskettedisp
{\vergleichskette
{ \delta(0)
}
{ =} { \gamma(0)
}
{ } {
}
{ } {
}
{ } {
}
}
{}{}{.}
Ferner ist

\mavergleichskettedisp
{\vergleichskette
{ \delta'( 0)
}
{ =} {  (1-x)  \begin{pmatrix}  0  \\ u     \end{pmatrix}
}
{ } {
}
{ } {
}
{ } {
}
}
{}{}{}
und

\mavergleichskettedisp
{\vergleichskette
{ \delta^{\prime \prime}( 0)
}
{ =} {  (1-x)  \begin{pmatrix}  -u^2  \\ 0     \end{pmatrix}
}
{ } {
}
{ } {
}
{ } {
}
}
{}{}{.}
Wenn die geforderten Bedingungen also erfüllt sein sollen, muss

\mavergleichskettedisp
{\vergleichskette
{ \omega
}
{ =} { (1-x) u
}
{ } {
}
{ } {
}
{ } {
}
}
{}{}{}
und

\mavergleichskettedisp
{\vergleichskette
{ \omega^2
}
{ =} { (1-x) u^2
}
{ } {
}
{ } {
}
{ } {
}
}
{}{}{}
gelten. Dies ergibt

\mavergleichskettedisp
{\vergleichskette
{ \omega^2
}
{ =} {  \omega u
}
{ } {
}
{ } {
}
{ } {
}
}
{}{}{}
und somit

\mavergleichskettedisp
{\vergleichskette
{ u
}
{ =} { \omega
}
{ } {
}
{ } {
}
{ } {
}
}
{}{}{}
und

\mavergleichskettedisp
{\vergleichskette
{ x
}
{ =} { 0
}
{ } {
}
{ } {
}
{ } {
}
}
{}{}{.}

 }

__NOEDITSECTION__

\inputaufgabeklausurloesung
{5}

{

Beweise den Satz über die Krümmung auf einer implizit gegebenen ebenen Kurve.

}
{

Es sei
\maabb {\gamma} { I } { \R^2
} {}
eine Bogenparametrisierung der Kurve $Y$ in einer Umgebung von $P$  mit

\mavergleichskette
{\vergleichskette
{ \gamma (0)
}
{ = }{ P
}
{  }{
}
{    }{
}
{   }{
}
}
{}{}{,}
die mit der gegebenen Orientierung übereinstimmt. Das totale Differential
\maabbdisp {\left(D N \right)_{ P }} { \R^2} {\R^2
} {}
ist linear, daher genügt es, die Aussage für den Vektor $\gamma'(0)$ zu zeigen. Nach
der [[Totale Differenzierbarkeit/K/Kettenregel/Fakt|Kettenregel]]
ist

\mavergleichskettedisp
{\vergleichskette
{ \left(  D N   \right)_{ P } \left(   \gamma'(0)  \right)
}
{ =} { (N \circ \gamma )'(0)
}
{ } {
}
{ } {
}
{ } {
}
}
{}{}{.}
Dieser Vektor ist ein Vielfaches von $\gamma'(0)$ und daher ist dies gleich

\mathdisp {\left\langle  (N \circ  \gamma)'(0) ,  \gamma' (0)   \right\rangle  \gamma'(0)} { . }

Wegen der Orthogonalitätsbedingung ist

\mavergleichskettedisp
{\vergleichskette
{ \left\langle  (N \circ  \gamma)(t) ,  \gamma' (t)   \right\rangle
}
{ =} { 0
}
{ } {
}
{ } {
}
{ } {
}
}
{}{}{}
für alle

\mavergleichskette
{\vergleichskette
{ t
}
{ \in }{ I
}
{  }{
}
{    }{
}
{   }{
}
}
{}{}{}
und daher

\mavergleichskettedisp
{\vergleichskette
{ \left\langle  (N \circ  \gamma)'(t) ,  \gamma' (t)   \right\rangle + \left\langle  (N \circ  \gamma)(t) ,  \gamma^{\prime \prime} (t)   \right\rangle
}
{ =} { 0
}
{ } {
}
{ } {
}
{ } {
}
}
{}{}{.}
Also ist

\mavergleichskettedisp
{\vergleichskette
{ \left(  D N   \right)_{ P } \left(   \gamma'(0)  \right)
}
{ =} { - \left\langle  (N \circ  \gamma)(0) ,  \gamma^{\prime \prime} (0)   \right\rangle  \gamma'(0)
}
{ =} { - \kappa( \gamma(0)) \gamma'(0)
}
{ } {
}
{ } {
}
}
{}{}{}
nach
[[Ebene reguläre Kurve/Krümmung/Skalarprodukt/Fakt|Lemma 3.8 (Differentialgeometrie (Osnabrück 2023))]].

 }

__NOEDITSECTION__

\inputaufgabeklausurloesung
{2}

{

Man gebe ein Beispiel für differenzierbare Mannigfaltigkeiten
\mathkor {} {M} {und} {N} {}
mit

\mavergleichskette
{\vergleichskette
{ \operatorname{ dim}_{  } ^{  }  \left(   M   \right), \, \operatorname{ dim}_{  } ^{  }  \left(   N   \right)
}
{ \geq }{ 1
}
{  }{
}
{    }{
}
{   }{
}
}
{}{}{}
und differenzierbare Abbildungen
\maabb {\varphi} { M } { N
} {}
und
\maabb {\psi} { N } { M
} {}
derart, dass

\mavergleichskette
{\vergleichskette
{ \psi \circ \varphi
}
{ = }{
\operatorname{Id}_{ M }
}
{  }{
}
{    }{
}
{   }{
}
}
{}{}{}
und

\mavergleichskette
{\vergleichskette
{\varphi \circ  \psi
}
{ \neq }{
\operatorname{Id}_{ N }
}
{  }{
}
{    }{
}
{   }{
}
}
{}{}{}
gilt.

}
{

Wir wählen

\mavergleichskette
{\vergleichskette
{M
}
{ = }{\R
}
{  }{
}
{    }{
}
{   }{
}
}
{}{}{}
und

\mavergleichskette
{\vergleichskette
{N
}
{ = }{\R^2
}
{  }{
}
{    }{
}
{   }{
}
}
{}{}{}
und
\maabbeledisp {\varphi} {\R} {\R^2
} { x } {(x,0)
} {,}
und
\maabbeledisp {\psi} {\R^2} {\R
} {(x,y)} { x
} {.}
Dann ist

\mavergleichskettedisp
{\vergleichskette
{ \left(  \psi \circ \varphi  \right) (x)
}
{ =} { \psi((x,0))
}
{ =} { x
}
{ } {
}
{ } {
}
}
{}{}{}
und

\mavergleichskettedisp
{\vergleichskette
{ \left(  \varphi  \circ \psi  \right) ((x,y))
}
{ =} { \varphi (x)
}
{ =} {(x,0)
}
{ } {
}
{ } {
}
}
{}{}{.}

 }

__NOEDITSECTION__

\inputaufgabeklausurloesung
{5}

{

Beweise die Aussage, dass die tangentiale Äquivalenz von Wegen auf einer Mannigfaltigkeit in einem Punkt $P$ mit einer beliebigen Karte überprüft werden kann.

}
{

Für eine
\definitionsverweis {differenzierbare Kurve}{}{}
\maabbdisp {\gamma} { I } { M
} {}
mit

\mavergleichskette
{\vergleichskette
{ 0
}
{ \in }{ I
}
{  }{
}
{    }{
}
{   }{
}
}
{}{}{}
und

\mavergleichskette
{\vergleichskette
{ \gamma(0)
}
{ = }{ P
}
{  }{
}
{    }{
}
{   }{
}
}
{}{}{}
und eine
\definitionsverweis {Karte}{}{}
\maabbdisp {\alpha} { U } { V
} {}
\zusatzklammer {mit

\mavergleichskettek
{\vergleichskettek
{P
}
{ \in }{ U
}
{  }{
}
{    }{
}
{   }{
}
}
{}{}{}
und

\mavergleichskette
{\vergleichskette
{V
}
{ \subseteq }{\R^n
}
{  }{
}
{    }{
}
{   }{
}
}
{}{}{}} {} {}
ändert sich der Ausdruck

\mathdisp {\left(  \alpha \circ \left(  \gamma \vert_{\gamma^{-1} (U)}  \right)  \right)'(0)} {  }

nicht, wenn man zu einem kleineren offenen Intervall

\mavergleichskette
{\vergleichskette
{ 0
}
{ \in }{ I'
}
{ \subseteq }{ I
}
{    }{
}
{   }{
}
}
{}{}{}
und einer kleineren offenen Menge

\mavergleichskette
{\vergleichskette
{P
}
{ \in }{ U'
}
{ \subseteq }{ U
}
{    }{
}
{   }{
}
}
{}{}{}
\zusatzklammer {mit der induzierten Karte} {} {}
übergeht. Wir können also davon ausgehen, dass
\mathkor {} {\gamma_1} {und} {\gamma_2} {}
auf dem gleichen Intervall definiert sind und ihre Bilder in $U$ liegen, und dass es für dieses $U$ zwei Karten
\maabbdisp {\alpha_1} { U } { V_1
} {}
und
\maabbdisp {\alpha_2} { U } { V_2
} {}
gibt. Dann folgt aus

\mavergleichskettedisp
{\vergleichskette
{ ( \alpha_1  \circ \gamma_1 )'(0)
}
{ =} { ( \alpha_1  \circ \gamma_2 )'(0)
}
{ } {
}
{ } {
}
{ } {
}
}
{}{}{}
nach
der [[Totale Differenzierbarkeit/K/Kettenregel/Fakt|Kettenregel]]
unter Verwendung der Differenzierbarkeit der
\definitionsverweis {Übergangsabbildung}{}{}

\mathl{\alpha_2 \circ \alpha_1^{-1}}{} sofort

\mavergleichskettealign
{\vergleichskettealign
{ \left(  \alpha_2 \circ \gamma_1  \right)'(0)
}
{ =} { \left(  \left(  \alpha_2 \circ \alpha_1^{-1}  \right) \circ \left(  \alpha_1 \circ \gamma_1  \right)  \right)'(0)
}
{ =} { \left(  D   \left(   \alpha_2 \circ \alpha_1^{-1}   \right)      \right)_{\alpha_1(P)} \left(   \left(  \alpha_1 \circ \gamma_1  \right)'(0)   \right)
}
{ =} { \left(  D   \left(   \alpha_2 \circ \alpha_1^{-1}   \right)      \right)_{\alpha_1(P)} \left(   \left(  \alpha_1 \circ \gamma_2  \right)'(0)   \right)
}
{ =} { \left(  \alpha_2 \circ \gamma_2  \right)'(0)
}
}
{}
{}{.}

 }

__NOEDITSECTION__

\inputaufgabeklausurloesung
{5}

{

Zeige, dass die Tangentialabbildung $T(\varphi)$ zu
\maabbeledisp {\varphi} {\R^1} {S^1
} { t } {( \cos  t, \sin  t  )
} {,}
surjektiv ist.

}
{

Die Abbildung $\varphi$ ist surjektiv, es ist also lediglich zu zeigen, dass für jedes
\mathl{t \in \R}{} die lineare Tangentialabbildung
\maabbdisp {} {T_t \R \cong \R} {   T_{\varphi(t)} S^1
} {}
surjektiv ist. Da beide Räume eindimensional sind, muss gezeigt werden, dass ein von $0$ verschiedener Vektor nicht auf $0$ geht. Ein Tangentialvektor an $t$ wird realisiert durch den differenzierbaren Weg
\maabbeledisp {\gamma} {\R} {\R
} { s } {t+s
} {.}
Der verknüpfte Weg
\maabbeledisp {\varphi \circ \gamma} {\R} {S^1
} { s } { ( \cos (t+s)  , \sin (t+s)  )
} {,}
realisiert den Bild-Tangentialvektor, und zwar ist
\zusatzklammer {in der umgebenden Ebene
\mathl{T_{\varphi(t)}S^1 \subseteq \R^2}{}} {} {}

\mavergleichskettedisp
{\vergleichskette
{( \varphi \circ \gamma)' (0)
}
{ =} {(- \sin  t , \cos  t  )
}
{ } {
}
{ } {
}
{ } {
}
}
{}{}{,}
und das ist nicht der Nullvektor.

 }

__NOEDITSECTION__

\inputaufgabeklausurloesung
{7 (3+4)}

{

Wir betrachten die differenzierbaren Abbildungen
\maabbeledisp {\gamma} { [1,c] } {\R^2
} { t } {(t,t^{3})
} {,}
\zusatzklammer {mit \mathlk{c \geq 1}{}} {} {}
und
\maabbeledisp {\varphi} {\R_+ \times \R_+ } {\R^3
} {(u,v)} {(u^3,u^2+v^2,u^{-1}v^{-1} )
} {,}
und die Differentialform

\mavergleichskettedisp
{\vergleichskette
{ \omega
}
{ =} { (x-y)dx-z^2dy+dz
}
{ } {
}
{ } {
}
{ } {
}
}
{}{}{}
auf dem $\R^3$.

a) Berechne die zurückgezogene Differentialform
\mathl{\varphi^* \omega}{} auf dem $\R_+ \times \R_+$.

b) Berechne das Wegintegral zur Differentialform
\mathl{\varphi^* \omega}{} zum Weg $\gamma$ in Abhängigkeit von $c$.

}
{

a) Die zurückgezogene Differentialform ist

\mavergleichskettealigndrucklinks
{\vergleichskettealigndrucklinks
{\varphi^* (  (x-y)dx-z^2dy+dz   )
}
{ =} { (u^3-u^2-v^2 ) d(u^3)  -u^{-2}v^{-2} d(u^2+v^2) +d (u^{-1}v^{-1})
}
{ =} { 3 (u^5-u^4-u^2v^2 ) du  -2u^{-1}v^{-2} du -2u^{-2}v^{-1} dv - u^{-2} v^{-1}du - u^{-1} v^{-2}dv
}
{ =} { (3 u^5-3u^4-3u^2v^2 -2u^{-1}v^{-2}- u^{-2} v^{-1} )  du + (  -2u^{-2}v^{-1} - u^{-1} v^{-2})dv
}
{ } {
}
}
{}{}{.}

b) Das Wegintegral ist

\mavergleichskettealigndrucklinks
{\vergleichskettealigndrucklinks
{ \int_1^c (3 t^5-3t^4-3t^2t^6 -2t^{-1}t^{-6}- t^{-2} t^{-3} )dt  +(  -2t^{-2}t^{-3} - t^{-1} t^{-6}) d t^3
}
{ =} { \int_1^c   (3 t^5-3t^4-3t^8-2t^{-7} - t^{-5})dt + ( -2t^{-5} - t^{-7})3t^2 dt
}
{ =} { \int_1^c (3 t^5-3t^4-3t^8-2t^{-7} - t^{-5}   -6t^{-3} - 3 t^{-5}      ) dt
}
{ =} { \int_1^c (-3t^8 + 3 t^5-3t^4 -6t^{-3} - 4 t^{-5}  -2t^{-7}       )  dt
}
{ =} { (-{ \frac{   1 }{  3 } } t^9 + { \frac{   1 }{  2 } } t^{6}  - { \frac{   3 }{  5 } } t^5  +  3 t^{-2} +  t^{-4} + { \frac{   1 }{  3 } }  t^{-6}     ) \vert_{  1 }^{  c  }
}
}
{
\vergleichskettefortsetzungalign
{ =} { -{ \frac{   1 }{  3 } } c^9 + { \frac{   1 }{  2 } } c^{6}  - { \frac{   3 }{  5 } } c^5  +  3c^{-2} +   c^{-4} + { \frac{   1 }{  3 } }  c^{-6}
-(-{ \frac{   1 }{  3 } } + { \frac{   1 }{  2 } } - { \frac{   3 }{  5 } } +3 + 1 + { \frac{   1 }{  3 } }  )
}
{ =} { -{ \frac{   1 }{  3 } } c^9 + { \frac{   1 }{  2 } } c^{6}  - { \frac{   3 }{  5 } } c^5 + 3 c^{-2} +  c^{-4} + { \frac{   1 }{  3 } } c^{-6} - { \frac{   39 }{  10 } }  }
{ } {}
{ } {}
}{}{.}

 }

__NOEDITSECTION__

\inputaufgabeklausurloesung
{3}

{

Es sei
\maabbdisp {\psi} { L } { M
} {}
eine
\definitionsverweis {differenzierbare Abbildung}{}{}
zwischen
\definitionsverweis {differenzierbaren Mannigfaltigkeiten}{}{,}
es sei
\maabbdisp {f} { M } {\R
} {}
eine Funktion auf $M$ und $\omega$ eine
$k$-\definitionsverweis {Form}{}{}
auf $M$. Zeige

\mavergleichskettedisp
{\vergleichskette
{ \psi^*(f \omega)
}
{ =} { \psi^*(f)  \psi^* \omega
}
{ } {
}
{ } {
}
{ } {
}
}
{}{}{.}

}
{

Es sei
\mathl{P \in L}{} und
\mathl{v_1 , \ldots , v_k \in T_PL}{.} Dann ist

\mavergleichskettealign
{\vergleichskettealign
{ \psi^* \left(  f \omega) (P,v_1 , \ldots , v_k  \right)
}
{ =} { (f \omega) \left(  \psi(P), T_P\psi (v_1) , \ldots , T_P\psi (v_k )       \right)
}
{ =} { f(\psi(P))   \cdot \omega \left(  \psi(P), T_P\psi (v_1) , \ldots , T_P\psi (v_k )       \right)
}
{ =} { \left(  \psi^* (f)  \right) (P)  \cdot \omega \left(  \psi(P), T_P\psi (v_1) , \ldots , T_P\psi (v_k )       \right)
}
{ =} { \left(  \psi^* (f)  \right) (P)   \cdot \left(  \psi^*( \omega)  \right)  \left(  P,v_1 , \ldots , v_k  \right)
}
}
{
\vergleichskettefortsetzungalign
{ =} { \left(  \psi^* (f)  \psi^*( \omega)  \right) \left(  P,v_1 , \ldots , v_k  \right)
}
{ } {}
{ } {}
{ } {}
}
{}{.}

 }

__NOEDITSECTION__

\inputaufgabeklausurloesung
{2}

{

Beschreibe möglichst viele
\definitionsverweis {zusammenhängende}{}{}
eindimensionale
\definitionsverweis {Mannigfaltigkeiten mit Rand}{}{.}

}
{  }

__NOEDITSECTION__

\inputaufgabeklausurloesung
{5}

{

Es sei
\maabbdisp {f} { [a,b] } { \R_{+}
} {}
eine
\definitionsverweis {stetig differenzierbare Funktion}{}{.}
Wir fassen den
\definitionsverweis {Subgraphen}{}{}
als eine
\definitionsverweis {Mannigfaltigkeit mit Rand}{}{}
auf, wobei der Rand aus dem Graphen, dem Grundintervall und den beiden Seitenkanten, aber ohne die vier Eckpunkte besteht. Bestätige den Satz von Stokes direkt für die
\definitionsverweis {Differentialform}{}{}

\mavergleichskette
{\vergleichskette
{ \omega
}
{ = }{ xdy
}
{  }{
}
{    }{
}
{   }{
}
}
{}{}{.}

}
{

Wegen

\mavergleichskettedisp
{\vergleichskette
{d(xdy)
}
{ =} { dx \wedge dy
}
{ } {
}
{ } {
}
{ } {
}
}
{}{}{}
ist einerseits

\mavergleichskettedisp
{\vergleichskette
{ \int_M dx \wedge dy
}
{ =} { \int_a^b f(t) dt
}
{ =} { F(b)-F(a)
}
{ } {
}
{ } {
}
}
{}{}{}
der Flächeninhalt des Subgraphen, wobei $F$ eine Stammfunktion von $f$ sei. Für das Kurvenintegral müssen wir den Rand von $M$ gegen den Uhrzeigersinn parametrisieren. Das Grundintervall wird durch
\mathl{i: t \mapsto \begin{pmatrix}  t  \\ 0   \end{pmatrix}}{} parametrisiert, dabei ist

\mavergleichskettedisp
{\vergleichskette
{ i^*\omega
}
{ =} { i^*xdy
}
{ =} { x d0
}
{ =} { 0
}
{ } {
}
}
{}{}{,}
dieser Ausschnitt des Wegintegrals ist also $0$. Für die Seite rechts ist
$i: t \mapsto \begin{pmatrix}  b  \\t   \end{pmatrix}$,
eine Parametrisierung, wobei $t$ zwischen
\mathkor {} {0} {und} {f(b)} {}
wandert. Dabei ist

\mavergleichskettedisp
{\vergleichskette
{ i^*xdy
}
{ =} { b dt
}
{ } {
}
{ } {
}
{ } {
}
}
{}{}{}
und das Wegintegral darüber ist
\mathl{b f(b)}{.} Der Beitrag der Kante links ist
\mathl{-af(a)}{.} Der Graph wird durch
\maabbeledisp {i} {[a,b]} {
} { t } {(t, f(t))
} {,}
parametrisiert. Dabei ist

\mavergleichskettedisp
{\vergleichskette
{ i^*xdy
}
{ =} { t df
}
{ =} { tf'(t) dt
}
{ } {
}
{ } {
}
}
{}{}{.}
Das Integral ist mit dem richtigen Vorzeichen gleich

\mavergleichskettealign
{\vergleichskettealign
{ - \int_a^b tf'(t)dt
}
{ =} { -  ( tf(t) - F  (t) ) \mid  _a^b
}
{ =} { - ( bf(b)-F(b)  -af(a) +F(a) )
}
{ =} { F(b)-F(a) +af(a) - bf(b)
}
{ } {
}
}
{}
{}{.}
Das gesamte Wegintegral ist somit ebenfalls

\mavergleichskettealign
{\vergleichskettealign
{bf(b) -  af(a)+   F(b)-F(a) +af(a) - bf(b)
}
{ =} { F(b)-F(a)
}
{ } {
}
{ } {
}
{ } {
}
}
{}
{}{.}

 }

__NOEDITSECTION__

\inputaufgabeklausurloesung
{9}

{

Beweise den Brouwerschen Fixpunktsatz.

}
{

Zur Notationsvereinfachung sei

\mavergleichskette
{\vergleichskette
{ r
}
{ = }{ 1
}
{  }{
}
{    }{
}
{   }{
}
}
{}{}{.}  Nehmen wir an, dass es eine fixpunktfreie stetig differenzierbare Abbildung $\psi$ geben würde. Dann ist stets

\mavergleichskettedisp
{\vergleichskette
{ x
}
{ \neq} { \psi(x)
}
{ } {
}
{ } {
}
{ } {
}
}
{}{}{,}
sodass die beiden Punkte eine Gerade definieren. Die Idee ist, mittels dieser Geraden einen
\zusatzklammer {der beiden} {} {}
Durchstoßungspunkt mit der Sphäre als Bildpunkt einer Retraktion auf den Rand zu nehmen. Mit der Hilfsfunktion

\mavergleichskettedisp
{\vergleichskette
{ h(x)
}
{ =} { { \frac{   \psi(x)-x }{  \Vert { \psi (x)-x} \Vert  } }
}
{ } {
}
{ } {
}
{ } {
}
}
{}{}{}
definieren wir eine Abbildung
\maabbdisp {\varphi} { B  \left(  0 , 1 \right) } { \R^n
} {}
durch

\mavergleichskettedisphandlinks
{\vergleichskettedisphandlinks
{ \varphi(x)
}
{ =} { x + \left(  - \left\langle  x  , h(x)  \right\rangle + \sqrt{1 + \left\langle  x  , h(x)  \right\rangle^2 - \Vert { x} \Vert^2  }  \right) \cdot h(x)
}
{ } {
}
{ } {
}
{ } {
}
}
{}{}{.}
Dabei ist der Ausdruck unter der Wurzel positiv. Dies ist bei

\mavergleichskette
{\vergleichskette
{ \Vert { x } \Vert
}
{ < }{ 1
}
{  }{
}
{    }{
}
{   }{
}
}
{}{}{}
klar und bei

\mavergleichskette
{\vergleichskette
{ \Vert { x } \Vert
}
{ = }{ 1
}
{  }{
}
{    }{
}
{   }{
}
}
{}{}{}
liegt ein Punkt auf der Sphäre vor, dessen Verbindungsgerade mit dem Kugelpunkt
\mathl{\psi(x)}{} nicht senkrecht zu $x$ ist
\zusatzklammer {der affine Tangentialraum zu einem Punkt der Sphäre trifft eine Kugel nur in einem Punkt} {} {,}
sodass

\mavergleichskettedisp
{\vergleichskette
{ \left\langle  x  , h(x)  \right\rangle
}
{ \neq} { 0
}
{ } {
}
{ } {
}
{ } {
}
}
{}{}{}
ist. Da die Quadratwurzel und der Betrag außerhalb des Nullpunktes
\definitionsverweis {stetig differenzierbar}{}{}
sind, handelt es sich bei
\mathkor {} {h(x)} {und bei} {\varphi(x)} {}
um stetig differenzierbare Abbildungen. Die Abbildung $\varphi$ bildet nach
[[Brouwersche Fixpunktsatz/Explizite Retraktion/Aufgabe|Aufgabe 23.23 (Differentialgeometrie (Osnabrück 2023))]]
die Kugel auf die Sphäre ab und ihre Einschränkung auf die Sphäre ist die Identität. Damit liegt eine stetig differenzierbare
\definitionsverweis {Retraktion}{}{}
der abgeschlossenen Vollkugel auf ihren Rand vor, was nach
[[Mannigfaltigkeit mit Rand/Stokes/Nichtexistenz von Retraktionen/Fakt|Satz 23.9 (Differentialgeometrie (Osnabrück 2023))]]
nicht sein kann.

 }

__NOEDITSECTION__

\inputaufgabeklausurloesung
{5 (3+1+1)}

{

Wir betrachten den Zylinder

\mavergleichskette
{\vergleichskette
{ Z
}
{ = }{ S^1 \times \R
}
{ \subseteq }{ \R^2 \times \R
}
{    }{
}
{   }{
}
}
{}{}{}
mit dem
\definitionsverweis {trivialen Zusammenhang}{}{}
auf seinem
\definitionsverweis {Tangentialbündel}{}{.}
\aufzaehlungdrei{Bestimme in einer geeigneten isometrischen Karte die
\definitionsverweis {geodätische Differentialgleichung}{}{}
für diejenige geodätische Kurve $\psi$  auf $M$,  die

\mavergleichskettedisp
{\vergleichskette
{ \psi(0)
}
{ =} { (1,0,3)
}
{ } {
}
{ } {
}
{ } {
}
}
{}{}{}
und

\mavergleichskettedisp
{\vergleichskette
{ \psi'(0)
}
{ =} { (0,1,-4)
}
{ } {
}
{ } {
}
{ } {
}
}
{}{}{}
erfüllt.
}{Löse die Differentialgleichung aus (1) in der Karte.
}{Finde eine geodätische Kurve in $Z$, die die Bedingungen aus (1) erfüllt.
}

}
{

\aufzaehlungdrei{Wir betrachten die isometrische
\zusatzklammer {inverse} {} {}
Karte
\maabbeledisp {\varphi} { ]- \pi, \pi[  \times \R } { S^1 \setminus \{  (-1,0) \} \times \R
} {(t,v)} { \left(  \cos  t   , \,   \sin  t  , \,  v                 \right)
} {.}
Dabei ist

\mavergleichskette
{\vergleichskette
{  \varphi( 0,3)
}
{ = }{ (1,0,3)
}
{  }{
}
{    }{
}
{   }{
}
}
{}{}{.}
Die Christoffelsymbole sind trivial, daher ist die geodätische Differentialgleichung gleich

\mavergleichskettedisp
{\vergleichskette
{ \gamma^{\prime \prime}
}
{ =} { 0
}
{ } {
}
{ } {
}
{ } {
}
}
{}{}{.}
Die angegebenen Bedingungen bedeuten auf der Karte

\mavergleichskettedisp
{\vergleichskette
{  \gamma(0)
}
{ =} { (0,3)
}
{ } {
}
{ } {
}
{ } {
}
}
{}{}{}
und

\mavergleichskettedisp
{\vergleichskette
{  \gamma'(0)
}
{ =} {  (1,-4)
}
{ } {
}
{ } {
}
{ } {
}
}
{}{}{.}
}{Eine Lösung in der Karte ist

\mavergleichskettedisp
{\vergleichskette
{  \gamma(t)
}
{ =} {  (0,3) + t (1,-4)
}
{ =} {  (t,3-4t)
}
{ } {
}
{ } {
}
}
{}{}{.}
}{Eine geodätische Kurve mit den Anfangsbedingungen ist daher

\mavergleichskettedisp
{\vergleichskette
{  \psi(t)
}
{ =} { \varphi( \gamma(t))
}
{ =} { \left(  \cos  t   , \,   \sin  t  , \,   3-4t                 \right)
}
{ } {
}
{ } {
}
}
{}{}{.}
}

 }

 [[Kategorie:Latexseite]]
